% =========================
% report/methodology.tex
% =========================

\subsection{فرمول‌بندی مسئله}
هدف یافتن هموگرافی $\mathbf{H}$ است که نقاط الگو را به نقاط متناظر در عکس نگاشت کند. اگر $\mathbf{x}_k^T$ نقطه‌ای در الگو و $\mathbf{x}_k^Q$ نقطهٔ متناظر در عکس باشد:
\[
\tilde{\mathbf{x}}_k^Q \sim \mathbf{H}\tilde{\mathbf{x}}_k^T.
\]
هموگرافی ۸ درجه آزادی مؤثر دارد (زیرا ضرب کل ماتریس در یک عدد غیرصفر همان نگاشت را می‌دهد). بنابراین برای تعیین $\mathbf{H}$ حداقل به ۴ تطبیق نقطه‌ای غیرهم‌خط نیاز داریم \cite{HartleyZisserman2004}.

\subsection{استخراج ویژگی با ORB}
ORB \cite{Rublee2011ORB} دو کار انجام می‌دهد:
\begin{itemize}
  \item \textbf{آشکارسازی نقاط کلیدی:} معمولاً بر مبنای FAST (گوشه‌ها).
  \item \textbf{توصیف‌گر دودویی:} بردار بیتی که با مقایسه شدت پیکسل‌ها در الگوهای مشخص ساخته می‌شود (BRIEF چرخش‌یافته).
\end{itemize}
مزیت توصیف‌گر دودویی این است که شباهت با فاصلهٔ همینگ (تعداد بیت‌های متفاوت) سنجیده می‌شود که سریع و مناسب برای کاربرد عملی است.

\subsection{تطبیق توصیف‌گرها}
فرض کنید برای هر تصویر مجموعه توصیف‌گرهای دودویی داریم. برای هر توصیف‌گر در الگو، نزدیک‌ترین همسایه در عکس را با فاصلهٔ همینگ پیدا می‌کنیم. سپس برای کاهش تطبیق‌های غلط:
\begin{itemize}
  \item \textbf{Cross-check:} فقط تطبیق‌هایی حفظ می‌شوند که دوطرفه نزدیک‌ترین همسایه باشند.
  \item \textbf{فیلتر فاصله:} تطبیق‌هایی که فاصلهٔ آنها از یک حد بیشتر است حذف می‌شوند (حد می‌تواند مطلق یا درصدی باشد).
\end{itemize}

\subsection{تخمین هموگرافی با RANSAC}
به دلیل وجود تطبیق‌های غلط، اگر مستقیماً $\mathbf{H}$ را با همهٔ تطبیق‌ها برازش کنیم، مدل منحرف می‌شود. RANSAC \cite{FischlerBolles1981} این مشکل را حل می‌کند:
\begin{enumerate}
  \item چند بار، ۴ تطبیق به‌صورت تصادفی انتخاب می‌شود و یک $\mathbf{H}$ کاندید ساخته می‌شود.
  \item برای هر تطبیق، خطای بازفرافکنی محاسبه می‌شود. اگر خطا کمتر از $\tau$ باشد، آن تطبیق \textbf{درون‌نهشتی} (Inlier) محسوب می‌شود.
  \item مدلی که بیشترین درون‌نهشتی را دارد انتخاب می‌شود و در نهایت روی مجموعهٔ درون‌نهشتی‌ها دوباره برازش می‌شود.
\end{enumerate}

\subsection{خطای بازفرافکنی (تعریف ریاضی و تفسیر)}
برای یک نقطهٔ الگو $\mathbf{x}^T=[x,y]^T$، ابتدا همگن می‌کنیم $\tilde{\mathbf{x}}^T=[x,y,1]^T$ و سپس:
\[
\tilde{\mathbf{y}} = \mathbf{H}\tilde{\mathbf{x}}^T = [u,v,w]^T,\quad
\mathbf{y} = \pi(\tilde{\mathbf{y}})=[u/w,\, v/w]^T.
\]
اگر $\mathbf{x}^Q$ نقطهٔ متناظر در عکس باشد، خطا:
\[
e = \|\mathbf{x}^Q - \mathbf{y}\|_2
\]
است که همان فاصلهٔ اقلیدسی (ریشهٔ جمع مربعات اختلاف مؤلفه‌ها) بر حسب پیکسل است. هرچه $e$ کوچک‌تر باشد، هم‌ترازی هندسی بهتر است.

\subsection{وارپ و هم‌پوشانی}
پس از تخمین $\mathbf{H}$، الگو به چارچوب عکس وارپ می‌شود (تبدیل پرسپکتیو). سپس خروجی هم‌پوشانی با آلفا-بلندینگ ساخته می‌شود:
\[
\mathcal{I}_{\text{out}} = \alpha \mathcal{I}_{\text{warp}} + (1-\alpha)\mathcal{I}_Q.
\]
در اینجا:
\begin{itemize}
  \item $\alpha$ بزرگ‌تر یعنی الگو پررنگ‌تر دیده می‌شود،
  \item $\alpha$ کوچک‌تر یعنی عکس واقعی غالب است.
\end{itemize}

\subsection{شاخص‌های اطمینان (Confidence Proxies)}
در کاربرد عملی، بهتر است یک شاخص اطمینان گزارش شود:
\begin{itemize}
  \item تعداد تطبیق‌های معتبر،
  \item نسبت درون‌نهشتی‌ها $N_i/N_m$،
  \item (در صورت زمین‌حقیقت) میانه/میانگین خطای بازفرافکنی.
\end{itemize}
این شاخص‌ها برای «پرچم‌گذاری» موارد مشکوک جهت بازبینی انسانی مفید هستند.
